% Fragmento de entrega (sin preambulo; se compila desde main.tex).
\begin{center}
\textbf{\large Memorándum para el consejo del banco: igualdad de trato y
costo de corregirla}
\end{center}

El modelo de riesgo que el banco ya usa rechaza créditos con un puntaje de
riesgo. Cuando se revisa por separado a hombres y mujeres, aparece una
desigualdad clara: entre los clientes que sí habrían pagado, se rechaza al
$36.74\%$ de los hombres y al $32.51\%$ de las mujeres. Esa diferencia de
$4.23$ puntos porcentuales se mantiene al repetir el análisis miles de veces
con muestras alternativas, por lo que no parece producto de la casualidad. En
cambio, la diferencia en el porcentaje de incumplidores que se aprueban
($25.03\%$ en hombres frente a $27.42\%$ en mujeres) sí podría deberse al
azar, y por eso no debería tratarse como una desigualdad probada. No es un
problema de intención: con probabilidades bien calibradas, un mismo corte
aplicado a todos produce resultados distintos cuando las tasas de
incumplimiento de cada grupo son distintas ($23.30\%$ en hombres y
$21.34\%$ en mujeres en la muestra de prueba).

La corrección que se propone es ajustar ligeramente el corte según el grupo
para igualar el porcentaje de rechazos entre quienes pagan. Con un corte de
$0.19$ para hombres y $0.18$ para mujeres, elegidos solo con datos de
entrenamiento, la diferencia en rechazos entre quienes pagan baja a $0.95$
puntos en la muestra de prueba, un valor ya indistinguible del azar. Las
demás brechas quedan también pequeñas ($0.78$ puntos en incumplidores
aprobados, $1.61$ en valor predictivo y $1.37$ en tasa de rechazo): con tasas
base de $23.3\%$ en hombres y $21.3\%$ en mujeres, la imposibilidad teórica
de igualar todas las medidas a la vez es leve en este caso concreto, y casi
todas las definiciones de trato justo se cumplen a la vez. Ojo: los intervalos
de confianza de esas brechas posteriores son anchos (entre $-0.012$ y $0.031$
para los rechazos de quienes pagan; entre $-0.034$ y $0.049$ para los
incumplidores aprobados), así que su compatibilidad con cero no prueba que
las brechas hayan desaparecido, solo que con esta muestra ya no se distinguen
del ruido. La ganancia pasa
de $1984$ a $1977$ unidades: la pérdida de $7$ unidades no se distingue de
cero (IC95\% entre $-100$ y $116$ unidades), así que no se puede afirmar que
corregir cueste dinero con esta muestra. Otras reglas (igualar la tasa de
rechazo o el valor predictivo) llegaron incluso a dar ganancias algo mayores
que la línea base ($+11$ y $+35$ unidades), también dentro del ruido, pero
dejan desigualdades de errores mayores y por eso no se recomiendan.

Hay compromisos que esta recomendación no elimina. Igualar un tipo de trato
no iguala todos los demás, aunque en este caso las demás diferencias quedan
pequeñas. Usar cortes distintos según el sexo es un trato diferenciado
explícito: necesita fundamento jurídico y de gobierno, no solo eficacia
técnica. Meter el sexo como dato del modelo tampoco resuelve nada: la
ganancia subiría unos $35$ unidades (IC95\% entre $-36$ y $100$, dentro del
ruido), pero la desigualdad entre grupos casi se duplica, y además otras
variables del expediente permiten adivinar el sexo con una exactitud
moderada. La recomendación práctica es, por tanto, mantener el modelo sin el
sexo como entrada, aplicar los cortes ajustados con revisión jurídica, y
vigilar periódicamente si la igualdad se sostiene.

